\section{七次射の有限性と標準偏極の豊富性}
\label{sec:septic-finiteness-v30}

前節の七次射は、一般点で有限であることまでしか用いなかった。
ここでは、同じ選択成分 $\mathcal Y_0$ の非空開部分上で、
曲面全体にわたり有限となることを証明する。
この区別は実質的である。単純分岐を検出した六変数部分族には、
実際に三本の収縮曲線がある。

\begin{theorem}[標準七次射の有限平坦性]
\label{thm:septic-finiteness-v30}
選択成分 $\mathcal Y_0$ のある非空開部分に対し、標準七次射
\[
 f:X=\operatorname{Bl}_{D}\mathbb P^2\longrightarrow\mathbb P^2
\]
は有限平坦で次数七である。したがって
\[
 L=5H-\sum_{i=1}^{18}E_i=f^*\mathcal O_{\mathbb P^2}(1)
\]
は豊富である。この性質は源と標的の射影座標変換で不変なので、
指定された七次元のPrym像の非空開部分にも成立する。
\end{theorem}

\subsection{収縮曲線の次数に対する一様な制限}

以下、既存の良い開部分上で $D$ は相異なる十八点、
$C$ は滑らかな六次曲線、$\eta$ はその二位数線束とする。
$D\sim K_C\otimes\eta$ および
\[
 h^0(C,\eta(H))=0,\qquad
 h^0(C,\eta(2H))=3
\]
であり、後者の完全線形系は基点を持たない。

\begin{lemma}[四次式に対するCayley--Bacharach性]
\label{lem:septic-CB4-v30}
$D$ の任意の十七点を通る四次曲線は存在しない。
\end{lemma}
\begin{proof}
$p\in D$ とする。平面上の制限完全列より
\[
 H^0(\mathbb P^2,\mathcal I_{D-p}(4))
 \simeq H^0(C,\eta(H+p)).
\]
$\deg\eta(H+p)=7$、$g(C)=10$ である。
Riemann--Rochと $\eta(2H)$ の基点自由性から
\[
 h^0(C,\eta(H+p))
 =h^0(C,\eta(2H-p))-2=2-2=0.
\]
\end{proof}

\begin{lemma}[収縮候補の有限な型]
\label{lem:septic-contraction-types-v30}
$f$ が収縮する既約曲線は、次のいずれかの平面曲線の狭義変換である。
\begin{enumerate}
\item $D$ の五点を通る直線。
\item $D$ の十点を通る二次曲線。
\item $D$ の十五点を通る三次曲線。
\item $D$ の一点で重複度二、別の十三点で重複度一を持つ三次曲線。
\end{enumerate}
\end{lemma}
\begin{proof}
$L\cdot E_i=1$ なので、例外曲線は収縮されない。
平面曲線の次数を $d>0$、十八点での重複度を $m_i\ge0$ と書く。
収縮条件は $\sum m_i=5d$ である。
その狭義変換の算術種数は非負なので、随伴公式から
\[
 \sum m_i^2\le d^2+2d+2.
\]
Cauchyの不等式と合わせると
\[
 \frac{25d^2}{18}\le d^2+2d+2,
 \qquad 7d^2-36d-36\le0,
\]
従って $d\le6$ である。
$d=5$ のとき、和が二十五である十八個の非負整数の平方和の最小値は三十九、
$d=6$ のときは五十四であり、上の上界三十七、五十にそれぞれ反する。
$d=4$ のとき
\[
 \sum_i m_i(m_i-1)\le6
\]
である。重複度が正の点の個数を $n$ とすると
$2(20-n)\le6$、従って $n\ge17$ となり、
補題\ref{lem:septic-CB4-v30}に反する。
$d=1,2,3$ について同じ平方和の上界を用いると、列挙した四つの型だけが残る。
\end{proof}

\subsection{六変数部分族の収縮曲線}

既存の同一成分の六変数部分族で
\[
 M=\begin{pmatrix}-1&a&c\\h&-1&j\\n&q&-1\end{pmatrix},\qquad
 (F_0,F_1,F_2)^t=M(x_0^2,x_1^2,x_2^2)^t
\]
と置く。$Q=\operatorname{diag}(F_0,F_1,F_2)+2xx^t$、$t=x$ であり、
\[
 (g_0,g_1,g_2)=(x_0F_1F_2,x_1F_0F_2,x_2F_0F_1).
\]
$M$ の成分、二次小行列式、行列式がすべて非零となる開部分を考える。
十八基点は十二個の $F_i=F_j=0$ と、六個の $x_i=F_i=0$ からなり、
すべて横断的である。

\begin{lemma}[三本の $(-6)$ 曲線]
\label{lem:septic-three-conics-v30}
この六変数部分族で、$f$ が収縮する既約曲線は
$F_0=0,F_1=0,F_2=0$ の狭義変換 $T_0,T_1,T_2$ に限る。
これらは互いに交わらない滑らかな有理曲線であり、
\[
 T_i^2=-6,\qquad L\cdot T_i=0.
\]
各 $T_i$ の像は標的の第 $i$ 座標点である。
\end{lemma}
\begin{proof}
各 $F_i$ は相異なる十基点を通る滑らかな二次曲線なので、
自己交叉数と収縮の主張は直ちに従う。
$F_i$ と $F_j$ の四交点はすべて一度ずつblow-upされるため、狭義変換は交わらない。
例外曲線は $L$ との交点数一なので収縮されない。

標的の三座標が非零の点 $[u_0:u_1:u_2]$ を取る。
$\prod x_iF_i\ne0$ の部分でファイバーの方程式を射影的に正規化すると
\[
 Y=\operatorname{diag}(u_0,u_1,u_2)M Y^{[2]}
\]
となる。この二次方程式系を斉次化したとき、無限遠の方程式は
$\operatorname{diag}(u)M Y^{[2]}=0$ であり、$\det M\ne0$ により解を持たない。
従ってこのアフィン方程式系は零次元である。

標的の一つの座標だけが零なら、$F_i=0$ を除いた源の逆像は
対応する座標直線 $x_i=0$ 上にある。その直線上の写像は、残りの添字を $j,k$ として
\[
 [x_j:x_k]\longmapsto
 [x_j(m_{kj}x_j^2-x_k^2):x_k(-x_j^2+m_{jk}x_k^2)]
\]
であり、仮定した非消滅条件により三次の非定数写像である。
標的が座標頂点なら、上記の三二次曲線の外での逆像は源の座標頂点に限られる。
以上から他の曲線は収縮されない。
\end{proof}

\subsection{同じ成分内での三本の非有効化}

前版で用いた有理点
\[
 (a,c,h,j,n,q)=(3,-1,4,-2,5/4,3/4),\qquad s=1
\]
を $A_*$ と書く。他の正規形パラメータは零である。
従属変数を $(p,m,i,r,o,l,b,g,f)$ とする九逆行列条件のJacobi行列式は
\[
 -\frac{4096}{60037250625}\ne0.
\]
従って $A_*$ は同じ $\mathcal Y_0$ の滑らかな点であり、
自由変数 $d,e,k$ の微分をそれぞれ一、他の自由変数の微分を零とする接ベクトルが存在する。
その非零の従属変数の微分は次のとおりである。
\[
\begin{array}{c|l}
 \text{自由方向}&\text{従属変数の微分}\\\hline
 d&\dot o=15/16,\quad\dot l=5/4,\quad\dot f=3/4\\
 e&\dot p=4,\quad\dot i=-8,\quad\dot g=-2\\
 k&\dot m=-1,\quad\dot r=3,\quad\dot b=-3
\end{array}
\]
これらは $\dot Z=-Z\dot A Z$ を九つの禁止成分で零とすることで計算される。

\begin{lemma}[三つの二次曲線条件への非接性]
\label{lem:septic-conic-obstruction-v30}
$A_*$ の近傍で十八基点を\'etale局所的に標識付ける。
$F_i=0$ 上にある十点が引き続き二次曲線上にあるという条件は、
$i=0,1,2$ のいずれについても $\mathcal Y_0$ の真の閉条件である。
\end{lemma}
\begin{proof}
十点の二次式評価行列を $E_i$ とする。$A_*$ では
$\ker E_i=\langle F_i\rangle$ である。
接ベクトルに沿う基点の速度を $\dot p_\nu$ とすれば、
二次式 $F_i+\epsilon H$ が動く十点を通るための必要十分な一次条件は
\[
 E_i H=b_i,\qquad
 (b_i)_\nu=-\nabla F_i(p_\nu)\cdot\dot p_\nu.
\]
基点の速度は $dg(p_\nu)\dot p_\nu+\dot g(p_\nu)=0$ を
射影アフィンチャートで解いて求める。

この線形系を $\mathbb F_{7561}$ 上で計算すると、十八基点はすべて分裂し、
各 $E_i$ の階数は五である。一方、付加行列の階数は
\[
\begin{array}{c|ccc}
 &d&e&k\\\hline
 F_0&6&6&5\\
 F_1&6&5&6\\
 F_2&5&6&6
\end{array}
\]
となる。従って $F_0,F_1$ には $d$ 方向、$F_2$ には $e$ 方向の一次延長が存在しない。

有限体計算から標数零への移行を明記する。
$A_*$、三接ベクトル、逆関数定理に使うJacobi小行列式の分母はすべて
$7561$ において単元である。十八基点はこの素数においても相異なり、
各基点で二つの局所方程式のJacobi行列は階数二である。
従って基点スキームは対応する局所整数環上で有限\'etaleである。
各十点への二次式評価写像も整数環上で定義され、
その特殊ファイバーは階数五、核は $F_i$ の直線である。
$5\times5$ 小行列式の一つが単元なので、二次式の五係数を五点の条件で消去できる。
残りの条件は底の正則関数となり、上の階数六はその微分の少なくとも一つが非零であることを示す。
これは標数零でも非零であり、当該閉条件は成分全体ではない。

\nolinkurl{septic_finiteness_v30.py} は、接ベクトル、十八点の座標、
速度、三つの評価行列、付加行列の非零小行列式を厳密に再計算する。
随伴JSONには行・列の添字と小行列式の値も保存する。
\end{proof}

\begin{proof}[定理\ref{thm:septic-finiteness-v30}の証明]
十八点を標識付ける有限\'etale被覆のうち、$A_*$ の上で選んだ点を通る既約成分 $S$ 上で考える。
補題\ref{lem:septic-contraction-types-v30}に現れる重複度型は有限個である。
各型について、指定された点でその重複度以上に消える次数 $d\le3$ の曲線が
存在するという条件は、平面曲線の射影係数空間との閉incidenceの固有像である。
従って閉条件である。

$A_*$ でそのような曲線が存在するなら、その全狭義変換と $L$ の交点数は零以下である。
$L$ はnefなので実際には零であり、全既約成分が収縮される。
補題\ref{lem:septic-three-conics-v30}から、それらは三つの二次曲線に限られる。
次数は三以下なので、候補は各 $F_i$ 自身のみである。
従って、$A_*$ を含む開近傍では、閉incidenceとして残る可能性はこの三つの十点二次曲線条件だけである。
補題\ref{lem:septic-conic-obstruction-v30}により、三条件はいずれも真の閉条件である。
既約な $S$ からこれらを除く非空開部分では、収縮曲線は存在しない。
この性質は標識付けに依存せず、\'etale写像は開写像なので、その像は $\mathcal Y_0$ の非空開部分を与える。

固有な全射 $f$ は曲線ファイバーを持たないので準有限、従って有限である。
滑らかな曲面から滑らかな曲面への有限射は、Cohen--Macaulay性と次元公式によって平坦である。
その次数は既に計算した $L^2=7$ である。
有限射による豊富な線束の引戻しは豊富なので、$L$ の豊富性も従う。
\end{proof}

\paragraph{この段階での分岐の数値.}
有限平坦な一般射でも、前節の
\[
 R_f\sim2\widetilde C,\qquad R_f^2=72,\qquad K_X\cdot R_f=0
\]
はそのまま成り立ち、分岐押出しサイクルの次数は二十四である。
分岐曲線の被約性と既約性は次節で証明し、分岐源の滑らかさは
定理\ref{thm:septic-ramification-smoothness-v30}で確定する。
